\section{P、NP 与 NP-hard：给问题按计算难度分类}
\label{sec:09-Convex-Optimization/09-Complexity-and-Approximation/01}

需求评审会上有两句话听起来是一回事。运营要``最优的投放组合''：预算固定，从上百个渠道里挑一批出来，让触达的人数最大。调度要``最优的派单方案''：把当天的订单分给司机，让总里程最小。两句话都是从一堆组合里挑一个，使某个量取到最好，而你手上是同一份求解器、同一台机器。

第二句可以答应。指派问题有多项式时间的精确算法，规模上到几万条订单也算得完，返回的方案带着``不可能更好''这句话。第一句不能。最大覆盖是 NP-hard 的，你能给出很好的组合，也能给出``最多比最优差多少''的保证，唯独不能保证它就是最优的那一个。这个差别与你的实现水平无关，也与有没有用上更快的库无关。

计算复杂度理论做的正是给需求分类这件事：把``难''从一句主观评价，变成问题自身的一项属性。

\subsection{多项式时间是一条画得很粗但站得很稳的线}

一个判定问题（答案只有是或否）属于 P，指的是存在确定性算法，在输入长度 \(n\) 的多项式时间内给出答案。排序、最短路、二分图匹配、线性规划都落在这一侧。

拿多项式而不是某个具体的阶数当分界线，看上去粗糙。\(O(n^{100})\) 的算法在 P 里，可它连一个像样的实例都跑不完。选这条线是因为它稳：换一种合理的机器模型、换一种合理的输入编码，一个算法是不是多项式时间不会改变，而``是否小于 \(n^3\)''会随着这些约定漂移。多项式类还对复合封闭——多项式次地调用另一个多项式算法，总量仍是多项式——于是子程序可以放心地拼起来，不必每拼一次重算一遍复杂度。

所以这条线划的是理论边界，不是工程上的``跑得动''。落在 P 里的问题未必跑得动；落在 P 外的问题也不是每个实例都跑不动。它的用处在更前一步：知道问题落在哪一侧，你才知道该不该继续去找那个多项式算法。

\subsection{验证容易，不保证求解容易}

NP 收的是另一种保证。一个问题属于 NP，指的是当答案为``是''时存在一份长度多项式的证书，拿着它可以在多项式时间内确认答案确实是``是''。哈密顿回路就是这样：有人递来一个顶点排列，你只需检查每个顶点恰好出现一次、相邻两点之间确有边。整数规划的可行性也是：有人递来一组整数赋值，你把它代进每条约束算一遍。两件事都不要求你说得出这份证书是怎么找到的。

\begin{center}
\begin{tikzpicture}[x=1cm,y=1cm,font=\small,>=stealth]
  \begin{scope}
    \draw[rounded corners=6pt,black!45] (0,0) rectangle (5.2,3.2);
    \node[font=\footnotesize,black!70] at (2.6,3.55) {候选解的集合：\(2^{n}\) 个子集};
    \foreach \x/\y in {0.6/0.7, 1.1/2.2, 1.7/1.2, 2.2/2.7, 2.6/0.5, 3.1/1.8,
                       3.6/2.9, 4.0/1.0, 4.6/2.3, 1.5/2.9, 3.4/0.9, 4.7/0.6,
                       0.8/1.6, 2.9/2.2, 2.0/1.9, 4.2/1.7, 1.3/0.4, 3.9/2.5}
      \fill[black!35] (\x,\y) circle (1.6pt);
    \fill[red!70!black] (3.1,1.8) circle (2.8pt);
    \node[font=\footnotesize,red!55!black] at (3.1,1.32) {\(x^\star\)};
    \draw[->,black!70] (-1.05,1.6) -- (-0.15,1.6);
    \node[font=\footnotesize,left,black!70] at (-1.1,1.6) {实例};
    \node[font=\footnotesize] at (2.6,-0.75) {求解：要在这里面定位那一个};
    \node[font=\footnotesize,black!55] at (2.6,-1.3) {没有已知的捷径};
  \end{scope}
  \begin{scope}[shift={(8.4,0)}]
    \node[draw,black!55,rounded corners=2pt,font=\footnotesize,
          minimum width=1.5cm,minimum height=0.55cm] (I) at (0.8,2.5) {实例};
    \node[draw,black!55,rounded corners=2pt,font=\footnotesize,
          minimum width=1.5cm,minimum height=0.55cm] (C) at (0.8,1.0) {证书};
    \node[draw,black!55,rounded corners=2pt,font=\footnotesize,
          minimum width=1.7cm,minimum height=1.1cm] (V) at (3.2,1.75) {检查器};
    \draw[->,black!70] (I) -- (V);
    \draw[->,black!70] (C) -- (V);
    \draw[->,black!70] (V) -- (5.0,1.75);
    \node[font=\footnotesize,right,black!70] at (5.05,1.75) {是};
    \node[font=\footnotesize] at (2.6,-0.75) {验证：把证书代进去核对};
    \node[font=\footnotesize,black!55] at (2.6,-1.3) {与读一遍输入同阶};
  \end{scope}
\end{tikzpicture}

\emph{同一个问题的两个方向。右边拿到一份候选答案，核对的成本与读一遍输入差不多；左边要在指数级的候选里定位那一个，至今没人知道怎样绕开这片集合。NP 说的是右边容易，P 说的是左边也容易。}
\end{center}

这幅图有一个现成的样本。\hyperref[sec:04-Discrete-Mathematics/03-Graph-Theory/07]{``欧拉回路与一笔画''}要求用尽每条边，判据是每个顶点的度数奇偶，判定与构造都在线性时间内完成；\hyperref[sec:04-Discrete-Mathematics/03-Graph-Theory/08]{``哈密顿回路与旅行商''}只把``边''换成``点''，问题就掉到图的左半边。两句话的表述几乎对称，成本相差一个量级。NP 这个类把那种不对称抽了出来，单独命名。

\(\text{P}\subseteq\text{NP}\) 一望即知：自己算得出答案，就不需要别人递证书。反向包含是否成立，就是 P 与 NP 是否相等这个问题，提出至今半个多世纪，仍然没有答案。普遍的猜测是两者不等，也就是存在验证得动、求解不动的问题。要留意，这只是猜测，而整个工程实践建立在这个猜测之上——密码学也一样。

还有一处容易滑手。NP 的定义只管``是''的那一侧。``这张图没有哈密顿回路''该怎样出示一份短的、可核对的证明，定义里没有要求，也没有人知道。把否定证书那一侧单独收起来的类叫 coNP。同时落在 NP 与 coNP 里的问题，例如线性规划的可行性和素数判定，两个方向都有短证书，一般认为它们比 NP 中最难的那批要容易些。

\subsection{归约把难度从一个问题搬到另一个}

多项式时间归约 \(A\le_p B\) 是一个多项式时间的变换，把 \(A\) 的实例映成 \(B\) 的实例，并保持答案的是与否。它可以朝两个方向读，读反了结论就反了。手上已经有 \(B\) 的快速算法时，\(A\le_p B\) 让你顺带解决 \(A\)；想说明 \(B\) 难时，要做的是从一个已知难的 \(A\) 归约过去。难度顺着箭头往前搬，不往回搬。

问题 \(B\) 称为 NP-hard，指的是 NP 中的每一个问题都能多项式时间归约到它。这个量词很强，它的含义是 \(B\) 至少和 NP 里最难的那些一样难：\(B\) 一旦有了多项式算法，整个 NP 会跟着塌进 P。NP-complete 是 NP-hard 且自己也落在 NP 里的那批问题，它们构成 NP 内部的最难层。Cook 与 Levin 在 1971 年证明可满足性问题 SAT 是 NP-complete，这是第一块砖；有了它，此后每证明一个新问题难，只需从已有的难题归约过去，不必再回到 NP 的定义。旅行商、整数规划、最大割、图着色、最大覆盖、集合覆盖都已经在这张名单上。具体的归约构造对使用者来说不必展开——你要做的是查表，不是重新证明一遍。

\subsection{NP-hard 要你放弃的是一个承诺，不是这件事本身}

判明需求落在 NP-hard 一侧，该改的是承诺的措辞。放弃的那句话是``总是给出最优解''，剩下的余地仍然很宽。

一条路是照样求精确解，只接受最坏情形上的指数。分支定界配上切平面，成熟的混合整数规划求解器每天在解上万变量的选址与排班问题。这类求解器能不能收敛，取决于每一步拿到的下界够不够紧，这一点在\hyperref[sec:09-Convex-Optimization/05-Duality/02]{``对偶间隙何时闭合''}里已经讲过：界紧，大片搜索空间当场被剪掉；界松，搜索树指数膨胀。最坏情形的指数是一句上界叙述，你手上的实例通常离最坏很远。

第二条路是近似算法：多项式时间内给出一个解，并证明它的目标值与最优值之比有最坏情形上限。度量旅行商的最小生成树法给出 2 倍保证，Christofides 算法把它压到 1.5 倍。这类保证的价值在于可以直接写进对外说明——``最多比最优差一半''是一句能被核对的话，而``效果不错''不是。

第三条是启发式：局部搜索、模拟退火、贪心加多次重启。它们在真实数据上常常打得过带保证的算法，却拿不出任何比例上限。风险不在平均表现，而在于你说不出哪一类实例会让它失手，也就没法把失手的代价提前算进方案里。

还有一条不显眼却常常最有效：利用实例的结构。NP-hard 讲的是最坏情形，而你拿到的实例总是有来历的。图是稀疏的，约束矩阵的系数只有 0 和 1 且排列得很规整，真正需要同时排的班次不超过几十个。参数化复杂度研究的就是``固定住某个参数之后变成多项式''这类情形，平滑分析则解释了单纯形法这样最坏指数的算法为何在实践中一路飞快。把实例的来历讲清楚，收益往往大过换一个更强的求解器。

\subsection{分类说的是问题，不是你手上的那个实例}

用这套语言时有几处边界值得先划清。

复杂度类定义在判定问题上，而工程里遇到的多是优化版本。``存在总里程不超过 \(B\) 的派单方案吗''属于 NP；``求最小总里程''一般不认为属于 NP，因为最优性本身没有短证书。所以对优化问题说 NP-hard 是准确的，说 NP-complete 通常不是。

近似算法有保证，不代表任意精度的近似都拿得到。除非 P = NP，一般的旅行商问题连常数倍的近似都不存在，度量情形下的 1.5 倍靠的是三角不等式这条额外结构。近似比也有它自己的难度理论，一个问题能被近似到什么程度，同样是问题的属性。

计算难度还不等于统计难度。一个问题可以算起来困难而所需数据很少，也可以算起来轻松而所需数据极多，两者各管各的。Part 13 的学习理论处理后者，本单元只处理前者。

分类到此为止能给的最实际的一句是：立项之前判断需求落在哪一边，比动手实现更要紧。落在 NP-hard 一侧之后，第一件要做的事是给自己弄一个可信的界——没有界，你连``现在这个方案差多少''都答不上来。而造界最常用的手法，是把令问题变难的那部分约束松开：让 0--1 变量在区间上连续取值，可行集立刻变成凸集，剩下的问题回到你会解的那一类。松开之后得到的当然不再是原问题的答案，但它是一个界，而且是一个能算出来的界。
